\section{Properties \& Limitations}

\begin{frame}{Desirable Properties of Curve Systems}
  \begin{columns}
    \begin{column}{0.5\textwidth}
      \begin{conceptbox}{What Makes a Good Curve?}
        \footnotesize
        \textbf{1. Intuitive Control:}
        \begin{itemize}
          \item Control points have geometric meaning
          \item Predictable curve behavior
          \item Easy to achieve desired shapes
        \end{itemize}

        \textbf{2. Local vs Global Control:}
        \begin{itemize}
          \item \textcolor{SecondaryColor}{Local}: Changing one control point affects only nearby curve
          \item \textcolor{AccentColor}{Global}: Changes affect entire curve
        \end{itemize}

        \textbf{3. Smoothness:}
        \begin{itemize}
          \item Continuous derivatives
          \item No unwanted oscillations
          \item Visually pleasing
        \end{itemize}

        \textbf{4. Computational Efficiency:}
        \begin{itemize}
          \item Fast evaluation
          \item Stable algorithms
          \item Memory efficient
        \end{itemize}
      \end{conceptbox}
    \end{column}
    \begin{column}{0.5\textwidth}
      \begin{center}
        \textbf{Local vs Global Control}

        \begin{tikzpicture}[scale=0.8]
          % Global control (Bézier)
          \begin{scope}[shift={(0,2)}]
            \coordinate (P0) at (0,0);
            \coordinate (P1) at (1,1);
            \coordinate (P2) at (2,0.5);
            \coordinate (P3) at (3,0);

            \draw[dashed, LightGray] (P0) -- (P1) -- (P2) -- (P3);
            \draw[thick, AccentColor] (P0) .. controls (P1) and (P2) .. (P3);

            \fill[AccentColor] (P0) circle (2pt);
            \fill[AccentColor!50] (P1) circle (3pt);
            \fill[AccentColor] (P2) circle (2pt);
            \fill[AccentColor] (P3) circle (2pt);

            \node[above] at (1.5,1.5) {\scriptsize \textcolor{AccentColor}{Global Control}};
            \node[below] at (1.5,-0.5) {\scriptsize Moving $P_1$ affects entire curve};
          \end{scope}

          % Local control (B-spline concept)
          \begin{scope}[shift={(0,0)}]
            \draw[thick, SecondaryColor] (0,0) .. controls (0.5,0.5) .. (1,0.3) .. controls (1.5,0.1) .. (2,0.5) .. controls (2.5,0.8) .. (3,0.5);

            % Show local influence
            \draw[dashed, SecondaryColor!50] (0.5,0) -- (0.5,1);
            \draw[dashed, SecondaryColor!50] (2,0) -- (2,1);

            \node[above] at (1.25,0.8) {\scriptsize \textcolor{SecondaryColor}{Local influence}};
            \node[below] at (1.5,-0.5) {\scriptsize Changes affect only nearby region};
          \end{scope}
        \end{tikzpicture}
      \end{center}
    \end{column}
  \end{columns}
\end{frame}

\begin{frame}{Bézier Curves: The Global Control Problem}
  \begin{columns}
    \begin{column}{0.6\textwidth}
      \begin{conceptbox}{Bézier Limitations}
        \footnotesize
        \textbf{1. Global Influence:}
        \begin{itemize}
          \item Every control point affects entire curve
          \item Cannot make local modifications
          \item Difficult to edit complex shapes
        \end{itemize}

        \textbf{2. High Degree Problems:}
        \begin{itemize}
          \item More control points = higher degree
          \item Potential for unwanted oscillations
          \item Numerical instability
        \end{itemize}

        \textbf{3. Limited Flexibility:}
        \begin{itemize}
          \item Fixed number of control points per segment
          \item Cannot represent some curves exactly (circles)
          \item Joining multiple segments requires care
        \end{itemize}
      \end{conceptbox}
    \end{column}
    \begin{column}{0.4\textwidth}
      \begin{center}
        \begin{tikzpicture}[scale=0.9]
          % High degree Bézier with unwanted oscillation
          \coordinate (P0) at (0,1);
          \coordinate (P1) at (0.5,2.5);
          \coordinate (P2) at (1,0.5);
          \coordinate (P3) at (1.5,2);
          \coordinate (P4) at (2,0.8);
          \coordinate (P5) at (2.5,1.8);
          \coordinate (P6) at (3,1.2);

          % Control polygon
          \draw[dashed, LightGray] (P0) -- (P1) -- (P2) -- (P3) -- (P4) -- (P5) -- (P6);

          % High degree curve (simulated)
          \draw[thick, AccentColor] (P0) .. controls (P1) and (P2) .. (1.2,1.3) .. controls (P3) and (P4) .. (2.2,1.1) .. controls (P5) and (P6) .. (P6);

          % Control points
          \foreach \p in {P0,P1,P2,P3,P4,P5,P6} {
            \fill[SecondaryColor] (\p) circle (1.5pt);
          }

          % Highlight problem area
          \draw[AccentColor, thick] (1.1,1.2) circle (0.2);
          \node[right] at (1.4,1.2) {\scriptsize \textcolor{AccentColor}{Unwanted oscillation}};

          \node[below] at (1.5,-0.5) {\scriptsize High degree = problems};
        \end{tikzpicture}
      \end{center}
    \end{column}
  \end{columns}
\end{frame}

\begin{frame}{Enter B-Splines: Local Control}
  \begin{columns}
    \begin{column}{0.5\textwidth}
      \begin{conceptbox}{B-Spline Advantages}
        \footnotesize
        \textbf{Local Control:}
        \begin{itemize}
          \item Each control point affects only nearby curve segment
          \item Easy local modifications
          \item Predictable editing behavior
        \end{itemize}

        \textbf{Flexible Degree:}
        \begin{itemize}
          \item Can have many control points with low degree
          \item No unwanted oscillations
          \item Numerically stable
        \end{itemize}

        \textbf{Continuity Control:}
        \begin{itemize}
          \item Automatic smoothness between segments
          \item Can control continuity level
          \item Professional modeling standard
        \end{itemize}
      \end{conceptbox}
    \end{column}
    \begin{column}{0.5\textwidth}
      \begin{center}
        \begin{tikzpicture}[scale=0.9]
          % B-spline with many control points
          \coordinate (Q0) at (0,1);
          \coordinate (Q1) at (0.5,2);
          \coordinate (Q2) at (1,0.5);
          \coordinate (Q3) at (1.5,1.5);
          \coordinate (Q4) at (2,1);
          \coordinate (Q5) at (2.5,2);
          \coordinate (Q6) at (3,1.2);

          % Control polygon
          \draw[dashed, LightGray] (Q0) -- (Q1) -- (Q2) -- (Q3) -- (Q4) -- (Q5) -- (Q6);

          % Smooth B-spline curve
          \draw[thick, SecondaryColor] (Q0) .. controls (Q1) .. (0.75,1.25) .. controls (Q2) .. (1.25,1) .. controls (Q3) .. (1.75,1.25) .. controls (Q4) .. (2.25,1.5) .. controls (Q5) .. (Q6);

          % Control points
          \foreach \q in {Q0,Q1,Q2,Q3,Q4,Q5,Q6} {
            \fill[SecondaryColor] (\q) circle (2pt);
          }

          % Show local influence
          \fill[SecondaryColor!20] (1.3,0.5) rectangle (1.7,2);
          \node[above] at (1.5,2.2) {\scriptsize Local influence};

          \node[below] at (1.5,-0.5) {\scriptsize Many points, low degree, smooth};
        \end{tikzpicture}
      \end{center}
    \end{column}
  \end{columns}
\end{frame}

\begin{frame}{When to Use What?}
  \begin{mathbox}{Curve Selection Guide}
    \footnotesize
    \begin{center}
      \begin{tabular}{|l|p{3cm}|p{3cm}|p{3cm}|}
        \hline
        \textbf{Application} & \textbf{Bézier Curves} & \textbf{B-Splines} & \textbf{NURBS} \\
        \hline
        \textbf{Simple shapes} & \textcolor{SecondaryColor}{\cmark} Good & \textcolor{LightGray}{-} Overkill & \textcolor{LightGray}{-} Overkill \\
        \hline
        \textbf{Typography} & \textcolor{SecondaryColor}{\cmark} Perfect & \textcolor{LightGray}{-} Too complex & \textcolor{LightGray}{-} Too complex \\
        \hline
        \textbf{UI Animation} & \textcolor{SecondaryColor}{\cmark} Ideal & \textcolor{LightGray}{-} Unnecessary & \textcolor{LightGray}{-} Unnecessary \\
        \hline
        \textbf{Complex curves} & \textcolor{AccentColor}{\xmark} Difficult & \textcolor{SecondaryColor}{\cmark} Excellent & \textcolor{SecondaryColor}{\cmark} Best \\
        \hline
        \textbf{CAD/Engineering} & \textcolor{AccentColor}{\xmark} Limited & \textcolor{SecondaryColor}{\cmark} Good & \textcolor{SecondaryColor}{\cmark} Standard \\
        \hline
        \textbf{3D Surfaces} & \textcolor{AccentColor}{\xmark} Patches only & \textcolor{SecondaryColor}{\cmark} Good & \textcolor{SecondaryColor}{\cmark} Perfect \\
        \hline
      \end{tabular}
    \end{center}
  \end{mathbox}

  \begin{conceptbox}{Recommendation}
    \footnotesize
    \textbf{Start with Bézier:} Simple, intuitive, widely supported \\
    \textbf{Move to B-splines:} When you need complex shapes with local control \\
    \textbf{Use NURBS:} For professional CAD and precise geometric modeling
  \end{conceptbox}
\end{frame}
